%% ============================================================
%%  Capítulo 6 – Requisitos del Fabricante de la PCB
%% ============================================================
\chapter{Requisitos del Fabricante de la PCB}
\label{chap:fabricante}

\section{Fabricante de Referencia}
\label{sec:fab_fabricante}

Se toma como fabricante de referencia \textbf{JLCPCB} (Shenzhen JLCPCB Co.,
Ltd.), seleccionado por los siguientes criterios:

\begin{itemize}
  \item Disponibilidad de proceso HDI con control de impedancia, necesario para
    la PCB-CTRL con BGA de 1056 pines.
  \item Documentación pública de design rules con valores numéricos concretos,
    que permite configurar las reglas directamente en Altium Designer.
  \item Costo competitivo para volúmenes de prototipo (1--10 unidades) y producción
    piloto (1\,000 unidades).
  \item Herramienta en línea de calculadora de impedancia controlada, utilizada
    para verificar el stack-up del Capítulo~\ref{sec:fab_stackup}.
\end{itemize}

\section{Capacidades Geométricas del Proceso}
\label{sec:fab_geometria}

Las restricciones geométricas que condicionan directamente el diseño de ALMA
son las siguientes. Estos valores han sido configurados como reglas DRC en el
proyecto de Altium Designer.

\begin{fabbox}[Reglas Geométricas del Fabricante -- JLCPCB]
  \begin{itemize}
    \item \textbf{Ancho mínimo de pista (track width):} 0.1\,mm (4\,mil).
      Determinante para el fan-out del BGA-1056 en la PCB-CTRL.
    \item \textbf{Separación mínima entre pistas (clearance):} 0.1\,mm (4\,mil).
    \item \textbf{Diámetro mínimo de perforación (drill hole):} 0.2\,mm
      (vías mecánicas con broca).
    \item \textbf{Anillo de cobre mínimo (annular ring):} 0.1\,mm.
    \item \textbf{Relación de aspecto de vías (aspect ratio):} máximo 8:1.
      Para un espesor de tarjeta de 1.6\,mm, el diámetro mínimo de perforación
      es $1.6/8 = 0.2$\,mm, consistente con el valor anterior.
    \item \textbf{Vías de paso (through-hole):} proceso estándar disponible en
      ambas PCBs.
    \item \textbf{Microvías (laser drill):} disponibles en proceso HDI para
      la PCB-CTRL (diámetro $\geq$0.1\,mm); no necesarias en PCB-PWR.
  \end{itemize}
\end{fabbox}

\begin{infobox}[Impacto en el diseño de la PCB-CTRL]
  El fan-out del BGA-1056 del A311D2 es el criterio geométrico más exigente del
  proyecto. Con un paso (pitch) de BGA de 0.65\,mm, el espacio disponible entre
  bolas es de 0.65\,$-$\,0.35 (diámetro bola)\,$=$\,0.30\,mm. Con el clearance
  mínimo de 0.1\,mm, es posible pasar \textbf{una pista de 0.1\,mm} entre dos
  bolas adyacentes del BGA, lo cual requiere disciplina máxima en el enrutamiento
  de la zona de escape.
\end{infobox}

\section{Stack-up y Número de Capas}
\label{sec:fab_stackup}

\subsection{Stack-up de PCB-CTRL}

El SoC A311D2 en BGA-1056 requiere un stack-up de alta densidad para:
(a)~rutar las señales de escape del BGA, (b)~proveer planos de referencia continuos
para MIPI DSI, y (c)~distribuir los múltiples rieles de alimentación dispersos en
el BGA.

\begin{fabbox}[Stack-up PCB-CTRL: 6 Capas]
  \begin{tabular}{cll}
    \toprule
    \textbf{Capa} & \textbf{Función} & \textbf{Tipo} \\
    \midrule
    1 (Top)    & Señales, componentes SMD       & Signal \\
    2          & Plano GND continuo             & Ground Plane \\
    3          & Señales internas (alta vel.)   & Signal \\
    4          & Plano de distribución potencia & Power Plane \\
    5          & Señales internas               & Signal \\
    6 (Bottom) & Señales, componentes SMD       & Signal \\
    \bottomrule
  \end{tabular}\\[6pt]
  \textbf{Espesor total:} 1.6\,mm (estándar) o 2.0\,mm (si se requiere mayor
  rigidez para el BGA).\\
  \textbf{Material:} FR-4, Dk~$\approx$~4.5 a 1\,GHz (para el cálculo de
  impedancias).
\end{fabbox}

\subsection{Stack-up de PCB-PWR}

La PCB-PWR no aloja señales de alta velocidad. Un stack-up de 4~capas es suficiente
y reduce el costo de fabricación.

\begin{fabbox}[Stack-up PCB-PWR: 4 Capas]
  \begin{tabular}{cll}
    \toprule
    \textbf{Capa} & \textbf{Función} & \textbf{Tipo} \\
    \midrule
    1 (Top)    & Componentes, pistas de potencia    & Signal/Power \\
    2          & Plano GND                          & Ground Plane \\
    3          & Plano de potencia (rieles DC)      & Power Plane \\
    4 (Bottom) & Pistas de retorno y pads térmicos  & Signal \\
    \bottomrule
  \end{tabular}\\[6pt]
  \textbf{Espesor total:} 1.6\,mm.
\end{fabbox}

\section{Control de Impedancia}
\label{sec:fab_impedancia}

JLCPCB ofrece control de impedancia con tolerancia de $\pm$10\,\%, suficiente para
los requisitos de ALMA.

\begin{fabbox}[Control de Impedancia -- PCB-CTRL]
  Las siguientes interfaces requieren control de impedancia activo en la PCB-CTRL:
  \begin{itemize}
    \item \textbf{MIPI DSI:} impedancia diferencial de 100\,$\Omega$~$\pm$~10\,\%
      en los 4~carriles (8~pistas en total). Topología microstrip diferencial en
      capa~1 con plano de referencia en capa~2 (GND).
    \item \textbf{USB 2.0:} impedancia diferencial de 90\,$\Omega$~$\pm$~10\,\%.
    \item \textbf{I2S / PDM:} rutas single-ended de 50\,$\Omega$ no requieren
      control de impedancia activo del fabricante, pero sí plano de referencia
      continuo.
  \end{itemize}
  La calculadora de impedancia del Layer Stack Manager de Altium Designer fue
  utilizada para determinar el ancho de pista~(W) y la separación~(S) de los pares
  diferenciales, verificados contra la calculadora en línea de JLCPCB.
\end{fabbox}

\section{Acabados Superficiales y Solder Mask}
\label{sec:fab_acabados}

\begin{itemize}
  \item \textbf{Acabado superficial:} ENIG (Electroless Nickel Immersion Gold)
    para ambas PCBs. Justificación: planitud superior para el montaje del BGA-1056
    y los componentes de paso fino (QFN, LGA) de los reguladores.
  \item \textbf{Solder mask:} verde, ambas caras. Apertura de solder mask sobre
    pads del BGA con expansión de $-$0.05\,mm (NSMD -- Non-Solder Mask Defined)
    según práctica recomendada para BGA de paso 0.65\,mm.
  \item \textbf{Serigrafía:} blanca, capa superior (mínimo para identificación de
    componentes y orientación).
\end{itemize}

\section{Perforación y Restricciones de Vías}
\label{sec:fab_vias}

\begin{itemize}
  \item \textbf{Vías de paso (through-hole):} diámetro de perforación mínimo
    0.2\,mm; se utiliza en ambas PCBs.
  \item \textbf{Vías térmicas (thermal vias):} en los pads de exposición térmica
    (EP) del amplificador Clase~D y los reguladores Buck con pad térmico. Diámetro
    0.3\,mm, rellenas o tapadas con solder mask para evitar fuga de soldadura
    durante el proceso de reflow.
  \item \textbf{Blind/buried vias:} no requeridas en esta iteración del diseño.
    El stack-up de 6~capas con through-hole vias es suficiente para el fan-out del
    BGA con el proceso de 0.1\,mm mínimo disponible.
\end{itemize}

\section{Configuración en Altium Designer}
\label{sec:fab_altium}

Las restricciones del fabricante han sido trasladadas a reglas DRC en el proyecto
de Altium Designer de la siguiente forma:

\begin{itemize}
  \item \textbf{Width Constraint:} mínimo 0.1\,mm configurado como regla global;
    regla específica de mayor ancho para las redes de potencia (VCC5, VDD\_CPU)
    calculadas con IPC-2152.
  \item \textbf{Clearance Constraint:} 0.1\,mm regla global; 0.5\,mm para redes
    de entrada de 12\,V respecto a redes de baja tensión (IEC\,62368-1).
  \item \textbf{Hole Size Constraint:} mínimo 0.2\,mm.
  \item \textbf{Annular Ring:} mínimo 0.1\,mm.
  \item \textbf{Impedance Profiles:} definidos en el Layer Stack Manager de Altium
    para los perfiles diferencial de 100\,$\Omega$ (MIPI) y 90\,$\Omega$ (USB).
  \item \textbf{Directiva SET PARAMETER:} se anotan en el esquemático de la
    PCB-CTRL los parámetros de fabricante relevantes (fabricante de referencia,
    número de capas, acabado superficial, control de impedancia requerido).
\end{itemize}
